\chapter{System Design}


\section{System Architecture}
SmartDock follows a classic three-tier, layered architecture (Figure~\ref{fig:arch}). The \textbf{presentation tier} is a single-page React application running in the user's browser. The \textbf{application tier} is a Node.js/Express REST API that also serves the built React files, so the whole system is one deployable web service. Requests pass through security middleware (rate limiting, JWT authentication, role and tenant checks) before reaching controllers, which delegate business logic to reusable services. The \textbf{data tier} is a MongoDB Atlas replica set, which provides the multi-document transactions required by the booking guarantee. Optional external providers deliver e-mail (Resend) and SMS (Twilio). Code pushed to GitHub is built and deployed automatically by Render.

\begin{figure}[!htbp]
\centering
\begin{tikzpicture}[node distance=0.35cm and 0.3cm]
  % users
  \node[blk, minimum width=3.1cm] (u1) {Warehouse / DC\\manager};
  \node[blk, minimum width=3.1cm, right=0.5cm of u1] (u2) {Carrier\\(transporter)};
  \node[blk, minimum width=3.1cm, right=0.5cm of u2] (u3) {Gate / security\\staff};
  % presentation
  \node[blk, minimum width=11cm, below=0.9cm of u2, fill=boxbg] (spa) {\textbf{React 19 single-page application} (Vite, Tailwind CSS, React Query)\\[1pt]\scriptsize Control tower $\cdot$ Carrier portal $\cdot$ Booking wizard $\cdot$ Analytics $\cdot$ Gate sheet $\cdot$ CSV/PDF export};
  \draw[arr] (u1) -- (u1 |- spa.north);
  \draw[arr] (u2) -- (spa.north);
  \draw[arr] (u3) -- (u3 |- spa.north);
  % api
  \node[blk, minimum width=11cm, below=0.9cm of spa, align=left, text width=10.6cm] (mw) {\textbf{Security middleware:} rate limiter $\rightarrow$ JWT authentication $\rightarrow$ role check (RBAC) $\rightarrow$ tenant / warehouse-access check};
  \node[blk, minimum width=11cm, below=0.25cm of mw, align=left, text width=10.6cm] (ctl) {\textbf{Controllers (REST):} auth, warehouses, docks, bookings, waitlist, vendors, analytics, audit-logs, notifications, demo};
  \node[blk, minimum width=11cm, below=0.25cm of ctl, align=left, text width=10.6cm] (svc) {\textbf{Services:} scheduling engine (slot grid, availability, atomic reservation) $\cdot$ waitlist $\cdot$ reliability scoring $\cdot$ notifications $\cdot$ audit $\cdot$ background jobs (no-show, expiry, reminders)};
  \node[blk, minimum width=11cm, below=0.25cm of svc, align=left, text width=10.6cm] (mdl) {\textbf{Data access:} Mongoose models with schema validation and indexes};
  \begin{scope}[on background layer]
    \coordinate (apitop) at ($(mw.north)+(0,0.42)$);
    \node[layer, fit=(apitop)(mw)(ctl)(svc)(mdl), label={[lbl, anchor=north west, xshift=3pt, yshift=-3pt]north west:Node.js + Express API (TypeScript) on Render}] (api) {};
  \end{scope}
  \draw[darr] (spa.south) -- node[left, font=\scriptsize]{HTTPS / JSON (REST)} (api.north -| spa.south);
  % db
  \node[blk, minimum width=6cm, below=1.0cm of api, fill=gray!8] (db) {\textbf{MongoDB Atlas} (replica set)\\\scriptsize users $\cdot$ warehouses $\cdot$ docks $\cdot$ bookings $\cdot$ vendoraccesses\\\scriptsize waitlistentries $\cdot$ auditlogs $\cdot$ notifications};
  \draw[darr] (api.south -| db.north) -- node[right, font=\scriptsize]{transactions} (db.north);
  % external
  \node[blk, right=0.4cm of db, minimum width=2.6cm, font=\scriptsize] (ext) {Resend (e-mail)\\Twilio (SMS)};
  \draw[arr] (api.south -| ext.north) -- (ext.north);
  \node[blk, left=0.4cm of db, minimum width=2.6cm, font=\scriptsize] (ci) {GitHub $\rightarrow$\\Render CI/CD};
  \draw[arr] (ci.north) -- (api.south -| ci.north);
\end{tikzpicture}

\caption{System architecture of SmartDock}
\label{fig:arch}
\end{figure}

\section{Data Flow Diagrams}
Figure~\ref{fig:dfd0} shows the context (level-0) diagram: SmartDock exchanges information with three human actors and the notification providers. Figure~\ref{fig:dfd1} decomposes the system into six processes and five data stores.

\begin{figure}[!htbp]
\centering
\begin{tikzpicture}
  \node[proc, minimum size=3cm, font=\small\bfseries] (sys) {0\\SmartDock\\system};
  \node[ent, left=3.2cm of sys] (car) {Carrier};
  \node[ent, right=3.2cm of sys] (mgr) {Warehouse\\manager};
  \node[ent, below=1.8cm of sys] (gate) {Gate staff};
  \node[ent, above=1.5cm of sys] (prov) {E-mail / SMS\\provider};
  \draw[flow] ([yshift=4mm]car.east) -- node[above]{access request, booking,\\reschedule, waitlist} ([yshift=4mm]sys.west);
  \draw[flow] ([yshift=-4mm]sys.west) -- node[below]{slots, status, score,\\notifications} ([yshift=-4mm]car.east);
  \draw[flow] ([yshift=4mm]mgr.west) -- node[above]{docks, rules, approvals,\\yard events} ([yshift=4mm]sys.east);
  \draw[flow] ([yshift=-4mm]sys.east) -- node[below]{schedule, KPIs,\\reports, audit} ([yshift=-4mm]mgr.west);
  \draw[flow] ([xshift=-5mm]gate.north) -- node[left]{check-in /\\departure} ([xshift=-5mm]sys.south);
  \draw[flow] ([xshift=5mm]sys.south) -- node[right]{gate sheet} ([xshift=5mm]gate.north);
  \draw[flow] (sys.north) -- node[right]{messages} (prov.south);
\end{tikzpicture}

\caption{Context diagram (DFD level 0)}
\label{fig:dfd0}
\end{figure}

\begin{figure}[!htbp]
\centering
\begin{tikzpicture}[x=1cm,y=1cm]
  \node[proc] (p1) at (0,0) {1.0\\Manage\\access};
  \node[proc] (p2) at (5,0) {2.0\\Configure\\site \& docks};
  \node[proc] (p3) at (10,0) {3.0\\Schedule\\appointments};
  \node[proc] (p4) at (10,-4.6) {4.0\\Execute\\yard};
  \node[proc] (p5) at (5,-4.6) {5.0\\Notify \&\\audit};
  \node[proc] (p6) at (0,-4.6) {6.0\\Analyse\\performance};
  \node[store] (d1) at (0,-2.3) {D1 Users / Access};
  \node[store] (d2) at (5,-2.3) {D2 Warehouses / Docks};
  \node[store] (d3) at (10,-2.3) {D3 Bookings};
  \node[store] (d4) at (13.3,-2.3) {D4 Waitlist};
  \node[store] (d5) at (5,-6.9) {D5 Audit / Notifications};
  \draw[flow] (p1) -- (d1);
  \draw[flow] (p2) -- (d2);
  \draw[flow] (d2) -- node[above, sloped]{rules, docks} (p3);
  \draw[flow] (d1) -- node[above, sloped, pos=0.35]{approved access} (p3);
  \draw[flow] (p3) -- (d3);
  \draw[flow] (p3.east) -| (d4.north);
  \draw[flow] (d4.south) |- node[below, pos=0.75]{auto-assign} (p4.east);
  \draw[flow] (d3) -- (p4);
  \draw[flow] (p4) -- node[above]{events} (p5);
  \draw[flow] (p5) -- (d5);
  \draw[flow] (d3.south west) -- node[above, sloped]{bookings, timestamps} (p6.north east);
\end{tikzpicture}

\caption{Level-1 data flow diagram}
\label{fig:dfd1}
\end{figure}

\section{Use Case Diagram}
Figure~\ref{fig:usecase} shows what each actor can do. The scheduler represents the automated background jobs.

\begin{figure}[!htbp]
\centering
\begin{tikzpicture}[scale=0.82, transform shape, uc/.style={ellipse, draw=brand, fill=boxbg, font=\scriptsize, minimum width=3.6cm, minimum height=0.75cm, align=center},
  actor/.style={font=\small\bfseries, align=center}]
  \node[draw=ink!50, rounded corners, minimum width=9.6cm, minimum height=12.4cm, label={[font=\small\bfseries]north:SmartDock}] (box) at (0,0) {};
  % carrier use cases (left column)
  \node[uc] (c1) at (-2.4,5.0) {Register \& request\\warehouse access};
  \node[uc] (c2) at (-2.4,3.5) {Book compatible slot};
  \node[uc] (c3) at (-2.4,2.0) {Book weekly series};
  \node[uc] (c4) at (-2.4,0.5) {Reschedule / cancel};
  \node[uc] (c5) at (-2.4,-1.0) {Join waitlist};
  \node[uc] (c6) at (-2.4,-2.5) {View score \&\\notifications};
  % manager use cases (right column)
  \node[uc] (m1) at (2.4,5.0) {Configure docks,\\hours, holidays};
  \node[uc] (m2) at (2.4,3.5) {Approve carriers};
  \node[uc] (m3) at (2.4,2.0) {Approve / reject\\appointments};
  \node[uc] (m4) at (2.4,0.5) {Gate check-in,\\dock start, complete};
  \node[uc] (m5) at (2.4,-1.0) {View KPIs \&\\export reports};
  \node[uc] (m6) at (2.4,-2.5) {Review audit log};
  % system
  \node[uc] (s1) at (0,-4.3) {Mark no-shows, expire requests,\\send reminders, auto-assign waitlist};
  % stick-figure actors
  \foreach \x/\name in {-6.4/Carrier, 6.4/{Warehouse\\manager}} {
    \begin{scope}[shift={(\x,1.2)}]
      \draw[thick, ink] (0,0.75) circle (0.2);
      \draw[thick, ink] (0,0.55) -- (0,-0.05) (-0.35,0.35) -- (0.35,0.35) (0,-0.05) -- (-0.25,-0.5) (0,-0.05) -- (0.25,-0.5);
      \node[actor] at (0,-0.95) {\name};
    \end{scope}
  }
  \coordinate (car) at (-6.0,1.3);
  \coordinate (mgr) at (6.0,1.3);
  \node[actor, draw=ink!50, rounded corners, inner sep=4pt] (sys) at (0,-7.4) {Scheduler (system clock)};
  \foreach \u in {c1,c2,c3,c4,c5,c6} \draw (car.east) -- (\u.west);
  \foreach \u in {m1,m2,m3,m4,m5,m6} \draw (mgr.west) -- (\u.east);
  \draw (sys.north) -- (s1.south);
\end{tikzpicture}

\caption{Use case diagram}
\label{fig:usecase}
\end{figure}

\section{Entity--Relationship Diagram}
Figure~\ref{fig:er} shows the main collections and their relationships. MongoDB stores them as documents; references are ObjectIds and are enforced in application code.

\begin{figure}[!htbp]
\centering
\begin{tikzpicture}[e/.style={rectangle split, rectangle split parts=2, draw=ink!80, fill=white, font=\scriptsize, align=left, rectangle split part fill={boxbg,white}, minimum width=3.4cm, inner sep=3pt},
  rel/.style={draw=ink!70, thick}, card/.style={font=\scriptsize\bfseries, color=brand}]
  \node[e] (user) at (0,0) {\textbf{USER}\nodepart{two}\_id (PK)\\username, email\\passwordHash, role\\warehouseId (FK)\\vendorCompany, status};
  \node[e] (wh) at (6,0) {\textbf{WAREHOUSE}\nodepart{two}\_id (PK), name, address\\timezone, operatingHours[7]\\slotDurationMinutes, bufferMinutes\\lateThreshold, noShowGrace\\blackoutDates[]};
  \node[e] (dock) at (12,0) {\textbf{DOCK}\nodepart{two}\_id (PK), warehouseId (FK)\\dockName, direction\\supportedTrailerTypes[]\\hazmatCertified, capacity\\maintenanceWindows[]};
  \node[e] (acc) at (3,-4.2) {\textbf{VENDOR\_ACCESS}\nodepart{two}vendorId (FK)\\warehouseId (FK)\\status, decidedBy};
  \node[e] (bk) at (9,-4.6) {\textbf{BOOKING}\nodepart{two}\_id (PK), warehouseId, dockId\\vendorId (FK), start/endTime\\status, direction, trailerType\\hazmat, PO/ASN, pallets\\checkedIn/dockStarted/completed\\arrivalStatus, lateBy\\recurrenceGroupId, history[]};
  \node[e] (wl) at (14,-4.2) {\textbf{WAITLIST\_ENTRY}\nodepart{two}vendorId, warehouseId\\dockId?, start/endTime\\load, truck, status\\fulfilledBookingId};
  \node[e] (au) at (3,-8.4) {\textbf{AUDIT\_LOG}\nodepart{two}warehouseId, actorId\\action, entity, message};
  \node[e] (nt) at (9,-9.4) {\textbf{NOTIFICATION}\nodepart{two}userId (FK), type, title\\message, read, email, sms};
  \draw[rel] (user) -- node[card, above]{1 \quad\quad\quad 1} (wh);
  \draw[rel] (wh) -- node[card, above]{1 \quad\quad\quad N} (dock);
  \draw[rel] (user) -- node[card, left, pos=0.3]{1}node[card, left, pos=0.85]{N} (acc);
  \draw[rel] (wh) -- node[card, right, pos=0.2]{1}node[card, right, pos=0.85]{N} (acc);
  \draw[rel] (dock) -- node[card, right, pos=0.15]{1}node[card, right, pos=0.85]{N} (bk);
  \draw[rel] (user.south east) -- node[card, above, pos=0.15]{1}node[card, above, pos=0.9]{N} (bk.west);
  \draw[rel] (wl) -- node[card, above, pos=0.2]{0..1}node[card, above, pos=0.85]{1} (bk);
  \draw[rel] (user.south) |- node[card, left, pos=0.3]{1 : N} (au.west);
  \draw[rel] (user.south) ++(-1,0) |- node[card, pos=0.9, above]{N} (nt.west);
\end{tikzpicture}

\caption{Entity--relationship diagram}
\label{fig:er}
\end{figure}

\section{Appointment Lifecycle}
Each booking follows a finite-state machine (Figure~\ref{fig:states}). Terminal states can never be revived, so capacity released by a cancellation cannot be silently re-occupied.

\begin{figure}[!htbp]
\centering
\begin{tikzpicture}[st/.style={rounded corners=8pt, draw=brand, fill=boxbg, font=\scriptsize\bfseries, minimum width=1.9cm, minimum height=0.8cm, align=center},
  term/.style={st, draw=ink!60, fill=gray!10}, t/.style={-{Stealth[length=2mm]}, draw=ink!70, font=\tiny, align=center}]
  \node[st] (p) at (0,0) {PENDING};
  \node[st] (a) at (3.4,0) {APPROVED};
  \node[st] (c) at (7.0,0) {CHECKED\_IN\\(in yard)};
  \node[st] (i) at (10.6,0) {IN\_PROGRESS\\(at dock)};
  \node[term] (d) at (14.3,0) {COMPLETED};
  \node[term] (r) at (0,-2.2) {REJECTED};
  \node[term] (x) at (5.2,-2.2) {CANCELLED};
  \node[term] (n) at (3.4,2.2) {NO\_SHOW};
  \draw[t] (p) -- node[above]{approve} (a);
  \draw[t] (a) -- node[below]{check-in} (c);
  \draw[t] (c) -- node[above]{start} (i);
  \draw[t] (i) -- node[below]{done} (d);
  \draw[t] (a.north east) to[bend left=25] node[above]{direct to door} (i.north west);
  \draw[t] (p) -- node[left]{reject} (r);
  \draw[t] (p) -- (x);
  \draw[t] (a) -- (x);
  \draw[t] (c) -- node[right]{cancel} (x);
  \draw[t] (a) -- node[right]{grace\\expired} (n);
\end{tikzpicture}

\caption{Booking state machine}
\label{fig:states}
\end{figure}

\section{Technology Stack}
\begin{table}[H]
\centering
\small
\caption{Technology stack}
\label{tab:stack}
\begin{tabularx}{\textwidth}{>{\bfseries\raggedright\arraybackslash}p{3.0cm}X}
\toprule
Layer & Technologies and purpose \\
\midrule
Language & TypeScript (strict mode) on both client and server \\
Frontend & React 19 \cite{react}, Vite 8, Tailwind CSS 4, TanStack React Query (data fetching, polling), React Router, React Hook Form + Zod (forms), date-fns, jsPDF + AutoTable (PDF export) \\
Backend & Node.js 22, Express 4 \cite{express}, Zod (request validation), JSON Web Tokens, bcrypt password hashing, express-rate-limit, Winston logging \\
Database & MongoDB Atlas (managed replica set) with Mongoose 8 ODM; multi-document transactions, compound indexes, aggregation pipelines \\
APIs & REST/JSON API (about 40 endpoints); Resend e-mail API and Twilio SMS API (optional) \\
Cloud \& DevOps & Render web service \cite{render} (Blueprint \texttt{render.yaml}, auto-deploy from GitHub), MongoDB Atlas free tier, GitHub for version control \\
Testing & End-to-end verification script with \texttt{mongodb-memory-server} (in-memory replica set) and two concurrent API instances; tenant-isolation unit test \\
Tools & Visual Studio Code, Git, npm, oxlint, Python + matplotlib (analysis and simulation), \LaTeX{} (this report) \\
\bottomrule
\end{tabularx}
\end{table}
